\documentclass[../main]{subfiles}
\begin{document}
\ifSubfilesClassLoaded{\setcounter{page}{4}}{}
\section{Inventing the “Bicycle Spectacles”}
\label{sec:01_inventing_spectacles}

What is communism? Where did it come from, where is it calling us, and whom is it calling? As History has shown, all these simple questions demand answers not once and for all, but every time they arise. As the screenwriter Grigori Gorin put it through the mouth of “The Very Same Munchhausen”, “love is a theorem that has to be proved every day”. It is the same for us, except that communism will be a theorem (a theory, even!) far more serious and far more global than love. And first we must learn to write down correctly, under the “given” of this theorem, everything that is genuinely present rather than merely imagined.

Objections to the very word “communism” are dictated by the historical regression our society has undergone. The vulgarisation of this concept after the collapse of the USSR, and even its transformation into a term of abuse, came about because of the triumph of social forces openly hostile to communism.

Even before the Soviet Union fell, the “blurring” of the concept had already begun, and so some theorists spoke of the need to draw distinctions within Soviet economics and ideology, to draw a dividing line so as not to mistake the inertial movement of capitalism for a communist tendency. Afterwards—because society had been plunged into the restoration of capitalism at a time when the tradition of communist theoretical thought had been all but lost, and the question was one of preserving at least something of it—the problems of communism and of the movement towards it were not specifically studied. All eyes were turned to flourishing capitalism. Both its supporters and its opponents spoke of it as something real, whereas communism and the movement towards it were seen as things of the past, or, in the case of its supporters, as a future so distant that one would first have to live long enough to see it, and so the important thing was to focus on surviving. For decades, communism as a theoretical problem ceased to exist.

Mass culture, far removed from theoretical thinking, has always actively reproduced (and still reproduces) the consciousness of the ruling class. This class fears and hates communism, clothing its fear and hatred in the form of various social myths on which a second generation is now growing up in the post-Soviet countries. Yet mythologised representations of communism have nothing in common with the concept of communism. Here we find the first crucial point of “hostility” (which is not actually hostility at all). Citizens of countries that slid into capitalism not because of foreign intervention but “under their own steam” simply no longer know what genuine, that is, scientific, communism is, or what the essence of the idea of communism is. And yet the problem of communism is once again confronting them as a practical one. And the integral theory whose trampling and neglect played no small part in the backslide into capitalism in the USSR and the other countries of the socialist camp needs to be revived and explained “in plain terms” to the people whose genuine interests it expresses. This applies not only to politically indifferent working proletarians, but also to those who consider themselves communists, or at least leftists—that is, to those who have politically defined themselves as sympathetic to communism. The latter perceive communism as some sort of “good” social order, interpreting this “good” rather at the level of sensory representations, often “by contradiction”—as the absence of the abominations and horrors of modern capitalist society. But such supporters cannot give a positive answer to the question “What is communism?” even to themselves.

This state of affairs is precisely what prompts us to produce this pamphlet now, as a pair of “bicycle spectacles” that allow one to see communism again and understand what it is. As for the opponents of communism, we must at the very least force them to raise not foolish but intelligent objections to communism—now scientific communism.

In the centenary year of October, interpretations of the role and conduct of the Bolsheviks in 1917 were so diverse in the public sphere that opinions and epithets predictably polarised: from noble heroes who “saved the state” to bandits who hated and destroyed that state (and statehood itself). But amid this polarisation the main point slipped away: the Bolsheviks, and of course Lenin, were thinkers, and in their actions they were guided not only by desires, preferences and so on, but also by a theory that was the culmination of the entire preceding development of human thought. The enemies of socialism and communism like to call October and the subsequent socialist construction an experiment performed by the Bolsheviks on the peoples… Well, these enemies are far closer to the truth than friends who are far removed from theory.

The Great October Revolution was precisely an experiment, provided, of course, that one does not understand the word foolishly, as some kind of external manipulation. Given how unprecedented it was, it could not have been anything else. It was an experiment conducted in keeping with every scientific tradition, for there is no experiment without a hypothesis. And there was such a hypothesis, developed by Lenin and his supporters and grounded in the legacy of Marx and Engels: the World Revolution. Having begun in Russia, the proletarian revolution was meant to stir up the industrially developed countries of Europe and set in motion their transformation into the World Republic of Soviets. By no means every scientific hypothesis is confirmed by experiment in its original formulation (even though the idea that proletarian revolutions are relevant remains broadly true today), yet even this experimental movement along a path charted by an even partly unconfirmed hypothesis moves science (socialism) forward.

The proposed metaphor simply brings us back to a different way of seeing the revolution, War Communism, the NEP, industrialisation and so on. It calls on us to view modern society and the trends in its development through the prism of communism as the most important science within the whole system of the social sciences that are needed in order to act rationally in today’s world. It is needed so that the revolutionary class can overcome its submission to constant exploitation, overcome within itself the oppressed class of capitalist society and take up its own class position, so that it can comprehend social relations and rationally change them towards the abolition of all oppression and of the exploitation of man by man, towards communism.

That is why the “bicycle spectacles” of the theory of communism, mocked by Mayakovsky in a different context, will only help us if we have the patience to “reinvent” them (to reproduce them): to travel Marx’s road to the Commune and the International, Lenin’s theoretical road to October, and then the entire Soviet road, without changing our “means of transport”—revolutionary theory.

To reinvent means to make the legacy of the science of communism our own under the present-day historical conditions we have inherited.
\end{document}
